\documentclass[10pt,a4paper]{amsart}
\usepackage[margin=19mm]{geometry}
\usepackage{amsmath,amssymb,mathtools}
\usepackage[scr=rsfs,cal=euler]{mathalfa}
\usepackage{xfrac}
\usepackage[unicode,hidelinks,bookmarks=false]{hyperref}
\newcommand{\ie}{i.e.\ }
\newcommand{\cf}{cf.\ }
\newcommand{\resp}{respectively\ }
\newcommand{\Subsection}{\subsection}
\providecommand{\nobreakdash}{\nobreak\hbox{-}}
\setlength{\emergencystretch}{2em}
\multlinegap0pt
\newtheorem{thm}{Theorem}
\newtheorem{lem}{Lemma}
\newtheorem{asp}{Assumption}
\newtheorem{defi}{Definition}
\newtheorem{rmk}{Remark}
\DeclareMathOperator{\ad}{ad}
\newenvironment{pf}{\noindent\textbf{Proof.}}{\hfill$\square$}
\title[Theorem 2]{Exponential time differencing for matrix-valued dynamical systems: global convergence of the METD$p$ scheme (Theorem 2)}
\author{Nayef Shkeir, Tobias Grafke}
\date{Source: arXiv:2406.13761v2 (CC BY 4.0)}
\begin{document}
\maketitle

\noindent\emph{Source and licence.} Exponential time differencing for matrix-valued dynamical systems, by Nayef Shkeir, Tobias Grafke. Version arXiv:2406.13761v2
(CC BY 4.0), distributed by arXiv under the Creative Commons Attribution 4.0 International licence,
http://creativecommons.org/licenses/by/4.0/, as recorded in the arXiv licence metadata for this exact
version and shown on the abstract page https://arxiv.org/abs/2406.13761v2. Full licence notice:
\texttt{licenses/arxiv-2406.13761-CC-BY-4.0.txt}.\par\smallskip

\noindent\textit{Editorial scope.} Counted unit: the article's Theorem 2 (source0.tex lines 457--473, source label \texttt{thm:metdp\_global}), the global convergence theorem for the arbitrary-order matrix exponential time differencing scheme METD$p$: under bounded generators $L,R$ with $[L,R]=0$, a globally Lipschitz nonlinearity and order-$p$ starting values, the global error on $[0,T]$ is bounded by $C\,h^{p}$. The article's complete original proof of Theorem 2 is delivered with it (source0.tex lines 475--612, one \texttt{pf} environment of 138 lines, adjacent to the statement, with no gap and no deferral). No mathematics is written, repaired or re-derived: the statement and the proof are the article's own lines, in the article's own notation (the update map $\Phi_h$, the one-step defect $\delta_{n+1}$, the local term $I_p$, the discrete error $E_n$), and no retained line is substituted, re-flowed or omitted, so nothing is removed and no omission record is needed. Retained uncounted as the source-local context the proof needs: the model problem and its nonlinear integral (lines 71--75, printed as equation (1), source label \texttt{eq:ODE}); the truncated local term $I_p$ and the METD$p$ update it defines (lines 271--282, printed as equations (19) and (20) in the article, source labels \texttt{eq:metdp\_Ip} and \texttt{eq:metdp\_scheme}); the article's Banach-algebra definition and its two standing assumptions (lines 329--338 and 345--348, printed as Definition 1, Assumption 1 and Assumption 2, source labels \texttt{def:1}, \texttt{asm:ass1} and \texttt{asm:ass2}); and the statement of the article's local defect theorem (lines 359--366, printed as Theorem 1, source label \texttt{thm:metdp\_local}), whose own proof is not retained, because the counted proof invokes that theorem as a black box. Every cross-reference of every retained line resolves inside this document. No retained line contains a citation, so the wrapper carries no bibliography and no bibliography entry. No retained line contains a figure environment, a graphics-inclusion command or drawing code, so no image is delivered and the package declares no asset dependency.

Editorial formatting only, in addition: an AMS \texttt{amsart} preamble that restates the article's own declarations and macros used by the retained lines -- the environments \texttt{thm}, \texttt{lem}, \texttt{asp}, \texttt{defi} and \texttt{rmk} with the article's own per-type, section-independent numbering; the operator declaration \texttt{\textbackslash ad}; and the article's own \texttt{pf} proof environment, whose definition \texttt{\textbackslash newenvironment\{pf\}\{\textbackslash noindent\textbackslash textbf\{Proof.\}\}\{\textbackslash hfill\$\textbackslash square\$\}} is copied from the article's own preamble (source0.tex line 31). Because the article numbers its theorem-like environments per type and not per section, the delivered numbering is exactly the article's own: Definition 1, Assumption 1, Assumption 2, Theorem 1 (the retained local defect theorem) and Theorem 2 (the counted theorem). Displayed equations are numbered consecutively in this document, so the five retained displayed items print here as equations (1)--(5) in order of appearance, whereas the article prints the same items as (1), (19), (20), (27) and (28).

The complete native source of the article as distributed in the arXiv e-print for this exact version is bundled unchanged as \texttt{sources/161045/source0.tex}: it is byte identical to the e-print file \texttt{main.tex} except that the pipeline's mechanical concatenation prepends one header comment line naming it. The article was typeset with the standard \texttt{article} class; no publisher class or style file is shipped in the e-print and none is redistributed or used.
Let us consider a matrix evolution equation for $Q(t) \in \mathbb{R}^{n \times n}$ given by 
\begin{equation}
    \label{eq:ODE}
    \dot Q = LQ + QR + N(Q,t),\quad Q(0) = Q_0\,,
\end{equation}

Define
\begin{equation}
\label{eq:metdp_Ip}
    I_p := h\sum_{m=0}^{p-1}\sum_{j=0}^{p-1-m} h^j\,C_{m,j}(A)\,
    \ad_R^{\,j}\!\left(\nabla^m N_n\right),
    \qquad A=h(L+R).
\end{equation}
The resulting order-$p$ METD multistep scheme is
\begin{equation}
\label{eq:metdp_scheme}
    Q_{n+1}=e^{hL}Q_ne^{hR}+I_p.
\end{equation}

\begin{defi}\label{def:1}
    Let $\mathcal{B}$ be a Banach algebra, i.e., a triple $(\mathcal{B}, \ast, ||\cdot||)$ where $(\mathcal{B}, \ast)$ is a unital associative algebra, and $(\mathcal{B},||\cdot||)$ is a (real or complex) Banach space, where the norm $||\cdot||$ is compatible with the multiplication $||x \ast y|| \leq ||x||\,||y||$ for any $x,y \in \mathcal{B}$. The commutator is then defined as $[x,y] = x \ast y - y \ast x$.
\end{defi}

In practice, we discretize our domain such that we will have $\mathcal{B}$ being the space of real (or complex) $n \times n$ matrices. We also make the following assumptions

\begin{asp}
    \label{asm:ass1}
    We assume that $L,R\in \mathcal B$ are bounded linear operators. Further, we assume that $L$ is an infinitesimal generator of an analytic semigroup $e^{tL}$ and the same for $R$ and $L+R$.
\end{asp}

\begin{asp}
    \label{asm:ass2}
    We assume that the nonlinearity $N : \mathcal{B} \times [0,T] \to \mathcal{B}$ is globally Lipschitz continuous on a strip along the exact solution. Moreover, along the exact solution the map $t\mapsto N(Q(t),t)$ is $C^{p}$ and its derivatives up to order $p$ are uniformly bounded on $[0,T]$.
\end{asp}

\begin{thm}
\label{thm:metdp_local}
Assume \ref{asm:ass1}--\ref{asm:ass2} and $[L,R]=0$. Then there exists a constant $C$ (independent of $h$ and $n$ with $t_{n+1}\le T$) such that the defect satisfies
\begin{equation*}
\|\delta_{n+1}\|\le C\,h^{p+1}.
\end{equation*}
In particular, METD$p$ is consistent of order $p$.
\end{thm}

\begin{thm}
\label{thm:metdp_global}
Assume \ref{asm:ass1}--\ref{asm:ass2} and $[L,R]=0$.
Let $Q(t)$ denote the exact solution of \eqref{eq:ODE} on $[0,T]$.
Let $\{Q_n\}_{n=0}^{M}$ be generated by the METD$p$ scheme \eqref{eq:metdp_scheme} with step size $h=T/M$.

Assume the starting values satisfy
\begin{equation}\label{eq:startup_assumption}
\max_{0\le n\le p-1}\|Q(t_n)-Q_n\|\le C_0\,h^p
\end{equation}
for some constant $C_0\ge 0$ (e.g.\ exact starting values, or produced by a one-step method of order $p$).

Then there exists a constant $C>0$, independent of $h$ (for $h\le h_0$) but depending on $T$, semigroup bounds for $e^{tL}$ and $e^{tR}$, the Lipschitz constant of $N$, and bounds on $t\mapsto N(Q(t),t)$ and its derivatives up to order $p$, such that
\begin{equation}\label{eq:global_error}
\max_{0\le n\le M}\|Q(t_n)-Q_n\|\le C\,h^{p}.
\end{equation}
\end{thm}

\begin{pf}
Let $\Phi_h:\mathcal{B}^p\to\mathcal{B}$ denote the METD$p$ update map, so that
\begin{equation*}
Q_{n+1}=\Phi_h(Q_n,\ldots,Q_{n-p+1})
      = e^{hL}Q_ne^{hR} + I_p(Q_n,\ldots,Q_{n-p+1}),
\end{equation*}
where $I_p$ is defined in \eqref{eq:metdp_Ip}. Define the global error
\begin{equation*}
e_n := Q(t_n)-Q_n,
\end{equation*}
and define the local defect by inserting the exact values into the update,
\begin{equation*}
\delta_{n+1}
:= Q(t_{n+1})-\Phi_h\bigl(Q(t_n),\ldots,Q(t_{n-p+1})\bigr).
\end{equation*}
By Theorem~\ref{thm:metdp_local}, there exists $C_{\mathrm{local}}>0$ such that
\begin{equation}\label{eq:glob_loc_def}
\|\delta_{n+1}\|\le C_{\mathrm{local}}\,h^{p+1}
\qquad\text{for all }n\text{ with }t_{n+1}\le T.
\end{equation}

Subtracting the numerical update from the exact update and using the definition of $\Phi_h$ we obtain the error recursion
\begin{equation}\label{eq:err_decomp}
e_{n+1}
= e^{hL}e_ne^{hR}
+ \Bigl(I_p(Q(t_n),\ldots,Q(t_{n-p+1}))-I_p(Q_n,\ldots,Q_{n-p+1})\Bigr)
+ \delta_{n+1}.
\end{equation}

Next, we claim that there exists a constant $K>0$, independent of $h$ and $n$ (for $t_{n+1}\le T$),
such that
\begin{equation}\label{eq:Ip_Lip}
\Bigl\|I_p(Q(t_n),\ldots,Q(t_{n-p+1}))-I_p(Q_n,\ldots,Q_{n-p+1})\Bigr\|
\le hK\,\max_{0\le \ell\le p-1}\|e_{n-\ell}\|.
\end{equation}
To see this, note from \eqref{eq:metdp_Ip} that $I_p$ is a finite sum (for fixed $p$) of terms of the form
\begin{equation*}
h^{1+j}\,C_{m,j}(A)\,\ad_R^{\,j}\bigl(\nabla^m N_n\bigr),
\qquad 0\le m\le p-1,\quad 0\le j\le p-1-m.
\end{equation*}
For $h\in(0,T]$, the operator coefficients $C_{m,j}(A)$ are bounded because they are finite linear
combinations of $\varphi$-functions evaluated at $A=h(L+R)$, and $\varphi_k$ are bounded operators under
Assumption~\ref{asm:ass1}. Moreover, $\ad_R$ is bounded since $R$ is bounded, hence $\ad_R^{\,j}$ is bounded for each
fixed $j$. Finally, by Assumption~\ref{asm:ass2}, $N(\cdot,t)$ is globally Lipschitz with constant $L_N$, so for each $\ell$
\begin{equation*}
\|N(Q(t_{n-\ell}),t_{n-\ell})-N(Q_{n-\ell},t_{n-\ell})\|\le L_N\|e_{n-\ell}\|.
\end{equation*}
Since $\nabla^m$ is a fixed finite linear combination of the past values $N_{n-\ell}$ (with coefficients depending only on $m$),
the difference $\nabla^m N(Q(t_\cdot))-\nabla^m N(Q_\cdot)$ is bounded by a constant (depending only on $p$) times
$\max_{0\le\ell\le p-1}\|e_{n-\ell}\|$. Collecting the finitely many bounded operator norms and combinatorial constants into $K$
yields \eqref{eq:Ip_Lip}.

Now, let us define 
\begin{equation*}
F_{n+1}
:= \Bigl(I_p(Q(t_n),\ldots,Q(t_{n-p+1}))-I_p(Q_n,\ldots,Q_{n-p+1})\Bigr)+\delta_{n+1}.
\end{equation*}
Then \eqref{eq:err_decomp} becomes the linear inhomogeneous recursion
\begin{equation}\label{eq:lin_inhom_rec}
e_{n+1}=e^{hL}e_ne^{hR}+F_{n+1}.
\end{equation}
Iterating \eqref{eq:lin_inhom_rec} gives the discrete variation-of-constants
\begin{equation*}
e_n
= e^{t_nL}e_0e^{t_nR}
 + \sum_{k=0}^{n-1} e^{(t_n-t_{k+1})L}\,F_{k+1}\,e^{(t_n-t_{k+1})R}.
\end{equation*}
Assuming exact initial data (or incorporating $e_0$ into the start-up assumption), we have $e_0=0$.
Using Assumption~\ref{asm:ass1}, for $0\le t_n-t_{k+1}\le T$,
\begin{equation*}
\|e^{(t_n-t_{k+1})L}\|\le C_L e^{\omega_L (t_n-t_{k+1})},\qquad
\|e^{(t_n-t_{k+1})R}\|\le C_R e^{\omega_R (t_n-t_{k+1})}.
\end{equation*}
Hence,
\begin{equation}\label{eq:en_sum_bound}
\|e_n\|
\le C_T\sum_{k=0}^{n-1}\|F_{k+1}\|,
\qquad
C_T:=C_LC_R\,e^{(\omega_L+\omega_R)T}.
\end{equation}
We define the maximum error
\begin{equation*}
   E_n:=\max_{0\le j\le n}\|e_j\|. 
\end{equation*}
Combining \eqref{eq:Ip_Lip} and \eqref{eq:glob_loc_def} yields, for all $k$ with $t_{k+1}\le T$,
\begin{equation*}
\|F_{k+1}\|\le hK\max_{0\le\ell\le p-1}\|e_{k-\ell}\| + C_{\mathrm{local}}h^{p+1}
\le hK\,E_k + C_{\mathrm{local}}h^{p+1}.
\end{equation*}
Inserting this into \eqref{eq:en_sum_bound} and using $nh\le T$ gives
\begin{align*}
\|e_n\|
&\le C_T\sum_{k=0}^{n-1}\bigl(hK E_k + C_{\mathrm{local}}h^{p+1}\bigr)
\le C_TK h\sum_{k=0}^{n-1}E_k + C_T C_{\mathrm{local}}\,T\,h^p.
\end{align*}
Taking the maximum over $0\le j\le n$ (so that $\|e_n\|\le E_n$) yields
\begin{equation}\label{eq:En_discrete}
E_n \le a h\sum_{k=0}^{n-1}E_k + b h^p,
\qquad
a:=C_TK,\quad b:=C_T C_{\mathrm{local}}\,T.
\end{equation}
For $n\ge p-1$, we can split the sum in \eqref{eq:En_discrete} as
\begin{equation*}
\sum_{k=0}^{n-1}E_k
\le \sum_{k=0}^{p-2}E_k + \sum_{k=p-1}^{n-1}E_k
\le (p-1)\max_{0\le j\le p-1}E_j + \sum_{k=p-1}^{n-1}E_k.
\end{equation*}
By the start-up assumption \eqref{eq:startup_assumption},
$\max_{0\le j\le p-1}E_j\le C_0 h^p$, so
\begin{equation*}
\sum_{k=0}^{n-1}E_k \le (p-1)C_0 h^p + \sum_{k=p-1}^{n-1}E_k.
\end{equation*}
Define $\widetilde E_n:=\max_{p-1\le j\le n}E_j$. Then from \eqref{eq:En_discrete} and the bound above,
for $n\ge p-1$,
\begin{equation*}
E_n \le a h\sum_{k=p-1}^{n-1}E_k + \bigl(a h (p-1)C_0 + b\bigr)h^p
\le a h\sum_{k=p-1}^{n-1}\widetilde E_k + \widetilde b\,h^p,
\end{equation*}
where $\widetilde b:=\bigl(a (p-1)C_0 + b\bigr)$.
Taking the maximum over $p-1\le j\le n$ gives
\begin{equation*}
\widetilde E_n \le a h\sum_{k=p-1}^{n-1}\widetilde E_k + \widetilde b\,h^p.
\end{equation*}
By the discrete Gr\"onwall inequality, this implies
\begin{equation*}
\widetilde E_n \le \widetilde b\,h^p e^{a(t_n-t_{p-1})}
\le \widetilde b\,e^{aT}\,h^p.
\end{equation*}
Since $E_n=\max\{\max_{0\le j\le p-1}E_j,\widetilde E_n\}$ and $\max_{0\le j\le p-1}E_j\le C_0 h^p$,
we conclude that for all $0\le n\le M$,
\begin{equation*}
E_n \le C\,h^p,
\qquad
C:=\max\{C_0,\widetilde b\,e^{aT}\}.
\end{equation*}
This proves \eqref{eq:global_error}.

\end{pf}

\end{document}
